\scanpage{637}
\mspage{\unsure{25}}

\uline{Théorème} 2.6. Soit $M$ une catégorie satisfaisant les
conditions suivantes

(2.6.1) \(\Bigg\{\)
a) $M$ est spéciale, \unsure{i.e.} stable par petites $\varinjlim$,
et ses éléments sont accessibles.

% slanted
\margin{\uline{NB} a $+$ c signifient que $M$ est une ``catégorie $U$-paratopos'', et impliquent l'existence d'une filtration cardinale satisfaisant à d) (sauf \unsure{l'assertion|l'accessibilité} \unsure{concernant les quotients} $Z'$ de $Z$) \unsure{cf.|c'est} \uline{\unsure{cat. de faisceaux !|de bien voir !}} Donc il s'agit d'un $U$-paratopos \unsure{satisfaisant en plus aux conditions} b) e) f) [\unsure{Est-on alors|Est-ce alors} \unsure{un hypo ?|une hypo ?} \uuline{Non}, \unsure{c'est|car} e), f)].}

b) \deleted{les limites \unsure{inductives}} filtrantes dans $M$ sont exactes à gauche \unsure{(}et dans $M$ les \deleted{\unsure{lim}} $\varprojlim$ finies sont représentables \unsure{(résulte de|résultent de} $a+c$)

c) $M$ admet une filtration cardinale
$(\mathrm{Filt}^{\pi})_{\pi\geq\pi_0}$.

d) les $\mathrm{Filt}^{\pi}$ sont stables par $\varinjlim$ indexées
par des cat. $I$ telles que $\operatorname{card}\mathrm{Fl}(I)\leq\pi$
et par quotients ($Z\in\mathrm{Filt}^{\pi}$ et $Z\longrightarrow Z'$
épi implique $Z'\in\mathrm{Filt}^{\pi}$)

\margin{\uline{NB} La condition d) est \unsure{sans doute inutile}, du moins si les épi sont effectifs, car d) sera satisfaite en tous cas pour $\pi=2^{\unsure{c|\ell}}$, avec $c\geq\pi_0$}

e) Si $A_1,A_2$ sont deux sous-objets d'un
objet $B$ de $M$, et $A_0=A_1\cap A_2$, alors
$A_1\amalg_{A_0}A_2\longrightarrow B$ est un mono.

% slanted
\margin{\uline{NB} La condition e) n'est \uline{pas} satisfaite pour la catégorie (Gr), et le th. est en défaut pour cette catégorie, laquelle n'admet \uline{pas} ``assez d'injectifs'' cf \deleted{\unsure{S}} SGA 4 I 8.17.9. Par contre, e) et f) sont satisfaites pour les catégories \uline{abéliennes}.}

f) Si on a un carré cocartésien

\diagram{\begin{tikzcd}[ampersand replacement=\&,row sep=small]
A_0 \arrow[r,"i"] \arrow[d] \& A_1 \arrow[d] \\
A_2 \arrow[r,"i'",no head] \& A_3
\end{tikzcd}}

avec $i$ un mono, alors $i'$ est mono.

% slanted
\margin{\uline{XVIII} cf Th. 6.7 (p. 157) pour une formulation plus \unsure{symétrique|synthétique} de l'ensemble des conditions a) à f) pour $M$.}

(On dira alors que $M$ est ``\uline{très spéciale}''.) Soit
de plus

\[
 W\subset\mathrm{Fl}\,M\tag*{(\unsure{2.6.2})}
\]

un localiseur de $M$, qui satisfait les conditions

(\unsure{2.6.3}) \(\Bigg\{\)
a) $W$ est stable par $\varinjlim$ filtrantes

% slanted
\margin{\uline{NB} a) implique déjà que $W=W^{\unsure{\natural|\sharp}}$}

b) $W$ est de plus hypo\uline{accessible}, i.e. il
existe une \uline{\unsure{petite}} partie $W_0\subset W$,
telle que $W$ soit la hypoenveloppe de
\unsure{$W$|$W_0$}, \inserted{pour $\unsure{\mathscr{V}|U}$ \unsure{sig.}} \unsure{i.e.} la plus petite partie de
$\mathrm{Fl}\,M$ \deleted{stable} contenant $W_0$, et stable
par $\varinjlim$ filtrantes.
% archive sheet number: 636
